\documentclass[12pt, a4paper]{article}
\usepackage[utf8]{inputenc}
\usepackage{amsmath, amssymb, amsthm, amsfonts,amsthm,tikz-cd}
\usepackage{marvosym}
\usepackage{mathrsfs}
\usepackage{CJKutf8}
\usepackage{bm}
\usepackage[margin=1in]{geometry}
\newcommand{\id}{\operatorname{id}}
\renewcommand{\hom}{\operatorname{Hom}}
\theoremstyle{remark}
\newtheorem*{remark}{Remark}
\newenvironment{lemma}[1]{
    \noindent \textbf{Lemma #1.} \itshape
}{
    \vspace{1em}
}

% 重新定义 solution 环境，保留标识并在末尾留出 8cm 空白
\newenvironment{solution}
  {\paragraph{\textit{Solution:}}\par\medskip}
  {\hfill$\blacksquare$\vspace{0.1cm}}

\title{Abstract Algebra III: Assignment - Week 5}
\author{
    \begin{CJK*}{UTF8}{gbsn}
    钟皓宇 (12533556)
    \end{CJK*}
}
\date{\today}

\begin{document}

\maketitle

\section*{Lemmas for Problem}
\begin{lemma}{A}
Let $\mu: N' \to N$ and $h: N' \to L$ be $R$-module homomorphisms. Let $Q = N \oplus_{N'} L$ be the pushout of $\mu$ along $h$, defined explicitly as the quotient module of the direct sum:
\[ Q = (N \oplus L) / W \]
where $W = \{(\mu(n'), -h(n')) \mid n' \in N'\}$. The induced maps are given by $\bar{h}: N \to Q, n \mapsto \overline{(n, 0)}$ and $i_L: L \to Q, l \mapsto \overline{(0, l)}$. Then the following properties hold:
\begin{enumerate}
    \item If $\mu$ is surjective, then the induced map $i_L$ is surjective.
    \item If $\mu$ is injective, then the induced map $i_L$ is injective.
\end{enumerate}
\end{lemma}
\begin{proof}
\textbf{(1)} Suppose $\mu$ is surjective. It suffices to prove that any pair with form $\overline{(n,0)}$ has preimage because $\overline{(n,l)} = \overline{(n,0)}+\overline{(0,l)}$ and $\overline{(0,l)}$ has preimage.
Because $\mu$ is surjective, then there exists $n'\in N'$ such that $\mu(n')=n$. I claim that $\overline{(n,0)} = \overline{(0,h(n'))}= i_L(h(n'))$ and thus surjectivity of $i_L$ holds.
To see that, we consider $\overline{(n,0)}-\overline{(0,h(n'))} = \overline{(n,-h(n'))}=0$ since $n=\mu(n')$.

\textbf{(2)} Suppose $\mu$ is injective. Let $l\in L$ be an arbitrary element such that $i_L(l)= \overline{(0,l)}=0$, or equivalently there exists $n'\in N'$ such that $\mu(n')=0$ and $-h(n')=l$.
The injectivity of $\mu$ implies $n'=0$ and hence $l=-h(n')=0$. Therefore, $i_L$ is injective.
\end{proof}

\begin{lemma}{B}
Let $f: M \to M''$ and $g: K \to M''$ be $R$-module homomorphisms. Let $P = M \times_{M''} K$ be the pullback of $f$ along $g$, defined explicitly by the subset of the direct sum:
\[ P = \{ (m, k) \in M \times K \mid f(m) = g(k) \} \]
equipped with the canonical projections $\pi_M: P \to M, (m, k) \mapsto m$ and $\pi_K: P \to K, (m, k) \mapsto k$. Then the following properties hold:
\begin{enumerate}
    \item If $f$ is surjective, then the projection $\pi_K$ is surjective.
    \item If $f$ is injective, then the projection $\pi_K$ is injective.
\end{enumerate}
\end{lemma}

\begin{proof}
\textbf{(1)} Suppose $f$ is surjective. Let $k \in K$ be an arbitrary element.
Because $f$ is surjective, then there exists an element $m\in M$ such that $f(m)=g(k)$, which implies $(m,k)\in P$ then $k=\pi_{K}(m,k)$ implies surjectivity.


\textbf{(2)} Suppose $f$ is injective. Suppose $(m,k)\in P$ such that $\pi_K(m,k)=0$. 
 Then we obtain that $f\circ \pi_M(m,k) = f(m) =g(k)= g\circ \pi_K(m,k) = 0$. The injectivity of $f$ implies $m=0$. Hence, $k=\pi_{K}(m,k)=0$ implies the injectivity of $\pi_K$.
\end{proof}

\section*{Problem 1}
Let $J$ be an $R$-module. Show that the following are equivalent.
\begin{enumerate}
    \item[(1)] $J$ is an injective module.
    \item[(2)] For any $R$-module monomorphism $\varphi : J \to M$, there exists an $R$-module homomorphism $\psi : M \to J$ such that $\psi\varphi = \mathrm{id}_J$.
    \item[(3)] If $J$ is a submodule of some $R$-module $M$, then $J$ is a direct summand of $M$.
\end{enumerate}
\begin{solution}
    We first prove that (1) implies (2). According to the definition of injective module, for morphism $\id_J$ and monomorphism $\varphi:J\to M$, there exists $\psi:M\to J$ such that diagram
    
    \begin{center}
        \begin{tikzcd}
J \arrow[r, "\varphi", hook] \arrow[dr, "\id_J"'] & M \arrow[d, "\exists \psi", dashed] \\
& J                                     
\end{tikzcd}
    \end{center}
commutes. In other words, $\psi\circ\varphi = \id_J$.

Next, we prove that (2) implies (3). Because $J$ is the submodule of $M$, then we obtain the canonical injection $\varphi:J\hookrightarrow M$, therefore (2) implies that there exists a morphism $\psi:M\to J$ such that $\psi\circ\varphi = \id_J$.
I claim that $M\cong J\oplus\ker(\psi)$. To see that, we define maps 
\begin{align*}
    f:J\oplus \ker(\psi) &\to M \\ (j,h) &\mapsto \varphi(j)+h 
\end{align*}
and 
\begin{align*}
        g:M &\to J\oplus \ker(\psi)  \\ m  &\mapsto (\psi(m),m-\varphi\circ\psi(m)).
\end{align*}
$f$ is a $R$-module morphism is quite clear because we have 
\begin{align*}
    \hom_R(J\oplus\ker(\varphi),M)\cong \hom_R(J,M) \oplus \hom_R(\ker(\varphi),M).
\end{align*}
We need to verify $g$ is a $R$-module morphism. The well-definedness of $g$ is given by the fact that 
\begin{align*}
    \psi(m-\varphi\circ\psi(m))=\psi(m) - \psi\circ \varphi\circ\psi(m)= 0,
\end{align*}
and $g$ is a $R$-module homomorphism since $\psi$ and are $R$-module homomorphism. I claim that $f,g$ are bijection and $f^{-1}=g$.
To see that, we just need to consider the compositon of two maps
\begin{align*}
    f\circ g(m)=\varphi\circ \psi(m)+(m-\varphi\circ \psi(m)) =m
\end{align*}
and 
\begin{align*}
    g\circ f(j,h) &= (\psi(\varphi(j)+h),\varphi(j)+h-\varphi\circ\psi(\varphi(j)+h)) \\
    &= (j+\psi(h),\varphi(j)+h-\varphi(j)+\varphi\circ\psi(h)) \\
    &=(j,h).
\end{align*}
Therefore, $M\cong J\oplus \ker(\psi)$ and then $J$ is a direct summand of $M$.

Now we prove (3) implies (1).  Suppose $\varphi:M\to N$ is an arbitrary monomorphism and $f:M\to J$ is an arbitrary homomorphism. 
Use the concepte already defined in \textbf{Lemma A.}
We consider the pushout $f$ along $\varphi$
\begin{center}
\begin{tikzcd}
M \arrow[r, "\varphi", hook] \arrow[d, "f"'] & N \arrow[d, "j"] \\
J \arrow[r, "i"', hook]  & J\oplus_{M}N
\end{tikzcd}
\end{center}
Since $\varphi$ is injective, $i$ is injective as well. Therefore, $J$ can be interpreted as the submodule of $J\oplus_{M}N$ by canonical inclusion. 
Hence, (3) implies that $J$ is the direct summand of $J\oplus_{M}N$, then we can assume that $J\oplus_{M}N\cong J\oplus K$. Define map $\psi=p\circ j:N\to J$ where 
$p:J\oplus K\to J$ is the canonical projection. Then $f=\psi\circ \varphi$ since $f=p\circ i\circ f = p\circ j\circ \varphi = \psi\circ \varphi$.
\end{solution}



\section*{Problem 2}
Suppose that $M_i$ ($i \in I$) are $R$-modules. Prove that $\bigoplus_{i \in I} M_i$ is a flat $R$-module if and only if each $M_i$, $i \in I$ is a flat $R$-module.
\begin{solution}
    Let $M=\bigoplus_{i\in I}M_i$. Let $f:A\hookrightarrow B$ be an arbitrary monomorphism. Assume that $M_i$ are all flat, then by definition, the $R$-module morphisms $f\otimes_R \id_{M_i}:A\otimes_R M_i\hookrightarrow B\otimes_R M_i$ are all
    monomorphic. Therefore
    \begin{align*}
        \bigoplus_{i\in I}(f\otimes_R \id_{M_i}):\bigoplus_{i\in I}(A\otimes_R M_i)\hookrightarrow \bigoplus_{i\in I} (B\otimes_R M_i)
    \end{align*}
    is monomorphism by definition of direct sum.
    Hence, the commutativity of tensor product and direct sum implies that 
    \begin{align*}
        f\otimes_R \id_{M}:A\otimes_R M \hookrightarrow B\otimes_R M
    \end{align*}
    is monomorphism meaning that $M$ is flat.

    Conversely, we use concepts defined in \textbf{Lemma B}. For arbitrary index $i$, consider canonical injection $j:B\otimes_R M_i\to B\otimes_R M$ pullback along $f\otimes_R \id_M$.
    Let $P= (B\otimes_R M_i)\times_{B\otimes_R M} (A\otimes_R M)$, and according \textbf{Lemma B}, the pullback homomorphism $g:P \to B\otimes_R M_i$ is monomorphism. Let $\iota:A\otimes_R M_i\to P$ by sending 
    pure tensor $a\otimes m_i$ to $(f(a)\otimes m_i,a\otimes m_i)$ and extending it linearly. It is clear that $\iota$ is well-defined because $f\otimes_R \id_M(a\otimes m_i) = f(a)\otimes m_i$ and for arbitrary $r\in R$, we have 
    \begin{align*}
        r.\iota(a\otimes m_i) = (f(r.a)\otimes m_i,r.a\otimes m_i) = \iota(r.a\otimes m_i)
    \end{align*}
    and other requirements of $R$-module is easy to prove. Therefore $\iota$ is a $R$-module morphism, hence it is injective since $a\otimes m_i$ equals zero if and only if $(f(a)\otimes m_i,a\otimes m_i)$ equals zero.
    Therefore $f\otimes_R \id_{M_i}= g\circ \iota$ is monomorphic meaning that $M_i$ is flat.
\end{solution}



\section*{Problem 3}
Prove: the following three conditions on an $R$-module $V$ are equivalent.
\begin{enumerate}
    \item[(1)] (Ascending chain condition.) For every infinite chain of $R$-submodules of $V$, 
    $$V_1 \subseteq V_2 \subseteq \cdots \subseteq V_n \subseteq \cdots,$$
    there is a number $m$ such that $V_m = V_{m+1} = \cdots$.
    \item[(2)] Any non-empty set of submodules of $V$ contains a maximal element with respect to inclusion.
    \item[(3)] Every $R$-submodule of $V$ is finitely generated.
\end{enumerate}
\begin{solution}
    First we show that (1) implies (2). If there exists a non-empty set of submodules of $V$, denoted by $\mathcal{S}$, such that no maximal element exists.
    Then for every submodule $V_1\in\mathcal{S}$, there exists a submodule $V_2\in\mathcal{S}$ such that $V_1\subsetneq V_2$. Continuing this process, we obtain an infinite ascending chain
    contradicting to (1).

    Now we prove (2) implies (3). Let $W$ be an arbitrary submodule of $V$. Let $\mathcal{S}$ be the collection of all finitely generated submodules of $W$. By (2), $\mathcal{S}$ contains a maximal element $W'$. Suppose $W'\subsetneq W$ is a proper submodule,
    then select an element $\alpha\in W\setminus W'$ and let $W''$ be the submodule generated by $W'$ and $\alpha$, which is finitely generated because $W'$ is finitely generated. However, $W'$ is the maximal element of $\mathcal{S}$
    meaning that $W''\notin \mathcal{S}$, which is a contradiction. Therefore, $W'=W$ implies that $W$ is finitely generated and thus
    every submodule of $V$ is finitely generated. 

    Finally, we claim that (3) implies (1). If there exists an infinite ascending chain 
    \begin{align*}
        V_1 \subsetneq V_2 \subsetneq \cdots \subsetneq V_n \subsetneq \cdots,
    \end{align*}
    By (3), the submodule $W= \cup_{i=1}^\infty V_i$ must be finitely generated.
    For arbitrary two elements $x,y\in W$, there exists an integer $n$ such that $x,y\in V_1\cup\dots\cup V_n= V_n$ by definition. Then all requirements of $R$-module are satisfied since $V_n$
    is a submodule. Suppose $W=\langle w_1, w_2, \dots, w_m \rangle$ is generated by finite elements $w_i$, then there exists an integer $s$ such that 
    $w_i\in V_1\cup\dots\cup V_s=V_s$ for all $i$ meaning that $W\subseteq V_s$ and therefore for every $k\ge s$ we obtain that $V_s=V_k$, which is a contradiction.
\end{solution}



\section*{Problem 4}
Let $W$ be an $R$-submodule of $V$. Prove: $V$ is Noetherian if and only if both $W$ and $V/W$ are Noetherian.
\begin{solution}
    Suppose $V$ is Noetherian. Every ascending chain of $W$ is also the ascending chain of $V$ meaning that chain is stationary. 
    Hence, Correspondence Theorem implies that if 
    \begin{align*}
        \overline{V}_1\subseteq \overline{V}_2\subseteq \dots  \tag{\Jupiter}
    \end{align*}
    is an ascending chain of $V/W$, there exists ascending chain of $V$
        \begin{align*}
        V_1\subseteq V_2\subseteq \dots \tag{\YinYang}
    \end{align*}
    such that $V_i$ is the lift of $\overline{V}_i$. The chain (\YinYang) is stationary implies chain (\Jupiter) is stationary because $\overline{V}_i = V_i+W$.

    Conversely, suppose $W$ and $V/W$ are Noetherian and let 
    \begin{align*}
        V_1\subseteq V_2\subseteq \dots \tag{\Gemini}
    \end{align*}
    be an arbitrary ascending chain of $V$. Since $V/W$ is Noetherian, the ascending chain 
    \begin{align*}
        V_1+W \subseteq V_2+W\subseteq \dots  \tag{\Scorpio}
    \end{align*}
    is stationary and we suppose it stationary at $n$. If chain (\Gemini) is not stationary, then select elements $\alpha_i\in V_i\setminus V_{i-1}$ for $i>n$.
    Chain (\Scorpio) stationary implies that $\alpha_i\in V_i+W = V_n+W$. Then we can suppose $\alpha_i = \beta_i+\gamma_i$ where $\beta_i\in V_n$ and $\gamma_i\in W$.
    Hence, $V_n\subseteq V_i$ implies that $\gamma_i\in V_{i}$ and $\gamma_i\notin V_{i-1}$. If not, $\alpha_i$ must be the element contained in $V_{i-1}$ contradicting our assumption.
    Consider ascending chain of $W$
    \begin{align*}
        \langle \gamma_{n+1} \rangle\subseteq \langle  \gamma_{n+1},\gamma_{n+2} \rangle \subseteq \dots
    \end{align*}
    which is stationary since $W$ is Noetherian. Suppose it is stationary at $m$ meaning that $\gamma_{m+1}$ can be linearly represented by $\gamma_{n+1},\dots,\gamma_{m}$
    and therefore $\gamma_{m+1}\in V_m$, which is a contradiction, chain (\Gemini) must be stationary.
\end{solution}


\section*{Problem 5}
Let $V$ be a Noetherian $R$-module and $f : V \to V$ an epimorphism. Show that $f$ is an $R$-isomorphism.

\begin{solution}
    Let $W_i=\ker(f^i)$ and consider an ascending chain
    \begin{align*}
        W_1\subseteq W_2\subseteq\dots \tag{\Gemini}
    \end{align*}
    because $V$ is Noetherian then this chain is stationary. Suppose it is stationary at $n$. 
    Then $W_n$ is the submodule of $V$ consisting exactly of all nilpotent elements associated to $f$.
    I claim that $f(W_n)\subseteq W_n$ and for every element $x\in V\setminus W_n$ we have $f(x)\notin W_n$.
    To see that, suppose $w\in W_n$ is nilpotent, then $f(w)$ is nilpotent as well, meaning that $f(w)\in W_n$.
    Hence, for arbitrary element $x\in V\setminus W_n$, if $f(x)\in W_n$ then it implies that  $f(x)$ is nilpotent and thus $x$ is nilpotent, which is contradiction.
    Suppose that $n>1$, ascending chain (\Gemini) stationary at $n$ means that $W_{n-1}\subsetneq W_{n}$, then $W_n\setminus W_{n-1}$ has no preimage of $f$ because $f(W_n)=W_{n-1}$.
    It is contradicts the assumption that $f$ is epimorphism.
    If $n=1$ and $W_1\neq \{0\}$ then $W_1 \setminus \{0\}$ has no preimage. Suppose $x\in V$ such that $f(x)\in W_1 \setminus \{0\}$, then $f^2(x) = 0$ implies that $x\in W_2$. 
    However, $W_2=W_1$ implies that $f(x)=0$, which is a contradiction. 
    Therefore, $f$ must be a monomorphism, hence an isomorphism. 
\end{solution}
\end{document}